%%%%%%%%%%%%%%%%%%%%%%
%%%%%%%%%%%%%%%%%%%%%%
\section{Annotation models}
%%%%%%%%%%%%%%%%%%%%%%
%%%%%%%%%%%%%%%%%%%%%%

We first collected all proteins with GO annotations supported by experimental evidence codes (EXP, IDA, IPI, IMP, IGI, IEP, TAS, IC) from the January 2011 version of the Swiss-Prot database (29,699 proteins in MFO; 31,608 in BPO; and 30,486 in CCO). We then generated three simple function annotation models: Na\"ive, BLAST and GOtcha, in order to assess the ability of performance metrics to accurately reflect the quality of a predicted set of annotations. In addition to these three methods, we generated another set of ``predictions" by collecting experimental annotations for the same set of proteins from a database generated by the GO Consortium released at about the same time as our version of Swiss-Prot. This was done to quantify the variability of experimental annotation across different databases using the same set of metrics. In addition, this comparison can be used to estimate the empirical upper limit of prediction accuracy because the observed performance is limited by the noise in experimental data. All computational methods were evaluated using 10-fold cross-validation. 

%%%%%%%%%%%%%%%%%%%%%%
%%%%%%%%%%%%%%%%%%%%%%
\subsection{The Na\"ive model}
\label{sec:naive}
%%%%%%%%%%%%%%%%%%%%%%
%%%%%%%%%%%%%%%%%%%%%%

The Na\"ive model was designed to reflect biases in the distribution of terms in the data set and was the simplest annotation model we employed. It was generated by first calculating the relative frequency of each term in the training data set. This value was then used as the prediction score for every protein in the test set; thus, every protein in the test partition was assigned an identical set of predictions over all functional terms. The performance of the Na\"ive model reflects what one could expect when annotating a protein with no knowledge about that protein.

%%%%%%%%%%%%%%%%%%%%%%
%%%%%%%%%%%%%%%%%%%%%%
\subsection{The BLAST model}
\label{sec:blast}
%%%%%%%%%%%%%%%%%%%%%%
%%%%%%%%%%%%%%%%%%%%%%

The BLAST model was generated using local sequence identity scores to annotate proteins. Given a target protein sequence $x$, a particular functional term $v$ in the ontology, and a set of sequences $S_{v}=\{s_{1},s_{2},\ldots\}$ annotated with term $v$, we determine the BLAST predictor score for function $v$ as 

%$\max\left\{ sid(x,s):s\in S_{v}\right\} $, 

\[
\underset{s\in S_{v}}{\max}\left\{ sid(x,s)\right\}, 
\]

\noindent where $sid(x,s)$ is the maximum sequence identity returned by the BLAST package \citep{Altschul1997} when the two sequences are aligned. We chose this method to mimic the performance one would expect if they simply used BLAST to transfer annotations between similar sequences.

%%%%%%%%%%%%%%%%%%%%%%
%%%%%%%%%%%%%%%%%%%%%%
\subsection{The GOtcha model}
%%%%%%%%%%%%%%%%%%%%%%
%%%%%%%%%%%%%%%%%%%%%%

The third method, GOtcha \citep{Martin2004}, was selected to incorporate not only sequence identity between protein sequences, but also the structure of the ontology (technically, BLAST also incorporates structure of the ontology but in a relatively trivial manner). %The GOtcha method outputs so-called i-scores for every protein by cumulatively propagating scores toward the root of the ontology and then normalizing them. As such it is an inexpensive and robust predictor of function. %
Specifically, given a target protein $x$, a particular functional term $v$, and a set of sequences $S_{v}=\{s_{1},s_{2},\ldots\}$ annotated with function $v$, one first determines the r-score for function $v$ as 

\[
r_{v}=c - \sum_{s\in S_{v}}\log(e(x,s)), 
\]

\noindent where $e(x,s)$ represents the E-value of the alignment between the target sequence $x$ and sequence $s$, and $c=2$ is a constant added to the given quantity to ensure all scores were above 0. Given the r-score for function $v$, i-scores were then calculated by dividing the r-score of each function by the score for the root term $i_{v}=r_{v}/r_{root}$. As such, GOtcha is an inexpensive and robust predictor of function.


